\documentclass{article}
% numbers — иначе natbib (автор-год) несовместим с plain thebibliography
\PassOptionsToPackage{numbers}{natbib}
\usepackage[preprint]{neurips_2025}
\usepackage{iftex}
\ifPDFTeX
  \usepackage[utf8]{inputenc}
  \usepackage[T1]{fontenc}
  \usepackage{times}
\else
  \usepackage{fontspec}
  % Latin Modern ships with tectonic; Liberation is not in the bundle.
  \setmainfont{lmroman10-regular.otf}
  \setmonofont{lmmono10-regular.otf}
\fi
\usepackage{graphicx}
\usepackage{booktabs}
\usepackage{amsmath,amssymb}
\usepackage{url}

\title{Diagnosing Search Ad Budget Waste: Segment Skew and Format--Device Mismatch in a Bicycle Retail Campaign}

\author{%
  KOI Research Agent\\
  ResearchOS\\
}

\begin{document}

\maketitle

\begin{abstract}
We analyze a May 2024 Yandex Direct search campaign for a regional bicycle retailer (Centra Market, Irkutsk) comprising 12{,}816 impressions, 1{,}391 clicks, and 50{,}005~RUB spend. Using a hypothesis-driven audit of campaign exports, we identify two structural drivers of inefficiency: (i)~budget concentration in low-click-through-rate (CTR) ad groups, and (ii)~a systematic format--device mismatch in which graphical ads on mobile achieve 3.32\% CTR versus 10.65\% for text ads on mobile---a 7.3 percentage-point gap that persists in all seven ad groups with paired formats. Counterfactual simulations on historical data show that reallocating 3{,}000~RUB from weak to strong groups yields only a modest +1.9\% click gain, whereas disabling mobile graphical ads and reinvesting freed budget into mobile text ads recovers nearly all lost clicks while raising campaign CTR from 10.85\% to 11.31\%. We discuss implications for practitioner-facing campaign diagnostics and the limits of observational counterfactuals.
\end{abstract}

\section{Introduction}

Paid search advertising allocates limited budgets across thousands of keyword--audience--creative combinations. Practitioners routinely observe spend flowing into segments with poor engagement, yet isolating \emph{why} budget is wasted---as opposed to merely \emph{where}---requires structured hypothesis testing rather than dashboard inspection alone.

This paper presents a case study of search campaign efficiency for Centra Market, a bicycle retailer in Irkutsk, Russia. We work from a single-month export (\texttt{bicycle.csv}, May 2024) and pursue two complementary hypotheses:

\begin{enumerate}
  \item \textbf{Segment skew:} A disproportionate share of budget flows to ad groups and query categories with below-median CTR, starving high-performing segments.
  \item \textbf{Format--device mismatch:} Graphical creatives underperform on mobile relative to text creatives, wasting spend independent of keyword targeting.
\end{enumerate}

For each hypothesis we specify measurable evidence criteria, run descriptive experiments on the export, and evaluate remediation policies via counterfactual budget simulations. Our contributions are methodological rather than algorithmic: we demonstrate how a small retail campaign can be diagnosed with transparent, reproducible analytics, and we report both positive and null-ish findings honestly---notably, group-level budget reallocation yields only modest gains, while format policy changes have a clearer efficiency signature.

\section{Related Work}

Search advertising optimization spans bid management \cite{edelman2007internet}, quality-score-aware budgeting \cite{ghosh2009ad}, and creative A/B testing \cite{blake2015consumer}. Device-specific performance heterogeneity is well documented in display and social advertising; mobile screen constraints and load latency plausibly amplify creative-format effects in search as well \cite{ghose2013mobile}. Our setting differs from large-scale auction simulations: we operate on a single advertiser's row-level export with no access to auction logs, conversion labels, or live experimentation infrastructure. The analysis therefore resembles \emph{marketing mix post-mortems} and \emph{observational policy evaluation} more than reinforcement-learning bid optimizers \cite{nuara2018combinatorial}.

\section{Data and Experimental Protocol}

\subsection{Dataset}

The campaign export contains 9{,}923 rows aggregated at the intersection of ad group, targeting category, ad format (text vs.\ graphical), and device type (mobile, desktop, tablet). Campaign-level totals for May 2024: 12{,}816 impressions, 1{,}391 clicks, 50{,}005~RUB spend, mean CTR $\approx$ 10.85\%, mean cost-per-click (CPC) 35.9~RUB. There are 25 named ad groups and six targeting categories.

\subsection{Metrics}

We use standard search metrics: CTR $= \text{clicks}/\text{impressions}$, CPC $= \text{spend}/\text{clicks}$, and clicks per 1{,}000~RUB spend (cl/1000RUB) as a budget-efficiency index. Statistical tests are descriptive; we do not claim causal identification beyond counterfactual accounting on fixed historical prices.

\subsection{Hypothesis Workflow}

Each hypothesis decomposes into \emph{cause evidence} (is the problem present?) and \emph{remediation} (does a plausible fix help on paper?). Experiments are tracked on a kanban board; we report only completed cards with filled result tables.

\section{Experiment 1: Segment-Level Budget Skew}

\subsection{Method}

We group rows by ad group name and compute impressions, clicks, CTR, spend, CPC, and budget share. We flag ``tail'' segments with CTR below 8\% and spend above 500~RUB. For remediation, we simulate transferring 1{,}000~RUB from each of three weak groups (\emph{Velosiped Katalog}, \emph{Velosiped Stels}, \emph{Velosiped Bolshoy}) to three strong groups (\emph{Velosiped Detskiy}, \emph{Velosiped Angarsk}, \emph{Velosiped Podrostkovyy}), holding total spend fixed and assuming marginal CPC in the recipient group.

\subsection{Results}

Table~\ref{tab:groups} contrasts exemplar high-spend groups. The top five groups by spend consume 60.3\% of budget at an average CTR of 9.4\%, below the campaign mean of 10.85\%. \emph{Velosiped Katalog} alone accounts for 9\% of spend (4{,}566~RUB) at 8.39\% CTR and 19.3 cl/1000RUB, while \emph{Velosiped Detskiy} spends 5.6\% of budget (2{,}780~RUB) at 22.13\% CTR and 29.9 cl/1000RUB.

\begin{table}[ht]
\centering
\caption{Selected ad groups illustrating budget--efficiency skew (May 2024).}
\label{tab:groups}
\begin{tabular}{lrrrr}
\toprule
Ad group & Spend (RUB) & CTR (\%) & Clicks & cl/1000RUB \\
\midrule
Velosiped Katalog & 4{,}566 & 8.39 & 87 & 19.3 \\
Velosiped Detskiy & 2{,}780 & 22.13 & 83 & 29.9 \\
Campaign average & -- & 10.85 & 1{,}391 & 27.8 \\
\bottomrule
\end{tabular}
\end{table}

The 3{,}000~RUB reallocation simulation increases total clicks from 1{,}391 to 1{,}418 (+26.7 clicks, +1.9\%). Pairwise transfers contribute +10.6 clicks (Katalog$\rightarrow$Detskiy), +10.6 (Stels$\rightarrow$Angarsk), and +5.5 (Bolshoy$\rightarrow$Podrostkovyy). Campaign-wide cl/1000RUB rises from 27.8 to 28.4.

\subsection{Interpretation}

Segment skew is \emph{confirmed}: budget concentrates in broad, lower-CTR groups while niche high-intent groups remain underfunded. However, remediation is \emph{moderate}: even a targeted 3{,}000~RUB shift---6\% of monthly spend---lifts clicks by less than 2\%. A backlog experiment (not executed) estimates that halving spend on broad/competitive/alternative query categories ($\sim$1{,}467~RUB) could yield $\sim$40 incremental clicks if reinvested at 36.6~RUB CPC; this remains prospective.

\section{Experiment 2: Format--Device Interaction}

\subsection{Method}

We cross-tabulate format and device over the full export, then repeat within ad groups where both text and graphical mobile creatives received $\geq$20 impressions. Success criteria: mobile graphical CTR at least 3 percentage points below mobile text CTR globally, and text winning in $\geq$5/7 paired groups.

For remediation we (a)~remove the mobile graphical segment (542 impressions, 18 clicks, 556~RUB) and recompute campaign totals, then (b)~reallocate the freed 556~RUB into mobile text at observed CPC 37.50~RUB.

\subsection{Results}

Table~\ref{tab:format} summarizes the four main format--device cells. Mobile graphical ads achieve 3.32\% CTR versus 10.65\% for mobile text ($\Delta = 7.3$~pp). Desktop graphical ads perform at 12.85\% CTR, comparable to the campaign average. Mobile text carries 9{,}229 impressions and 983 clicks (36{,}863~RUB).

\begin{table}[ht]
\centering
\caption{CTR by ad format and device type (May 2024).}
\label{tab:format}
\begin{tabular}{llrrrr}
\toprule
Format & Device & Impr. & Clicks & CTR (\%) & CPC (RUB) \\
\midrule
Text & Mobile & 9{,}229 & 983 & 10.65 & 37.50 \\
Graphical & Mobile & 542 & 18 & 3.32 & 30.87 \\
Text & Desktop & 14 & 4 & 28.57 & 34.50 \\
Graphical & Desktop & 2{,}826 & 363 & 12.85 & 32.59 \\
\bottomrule
\end{tabular}
\end{table}

In paired ad-group comparisons, text beats graphical on mobile in \textbf{7/7} groups (gaps +1.0 to +11.0~pp). Examples: \emph{Velosiped Kupit} 9.31\% vs.\ 2.70\%; \emph{Velosiped Angarsk} 14.35\% vs.\ 3.33\%.

Policy simulation results appear in Table~\ref{tab:policy}. Removing mobile graphical ads raises CTR from 10.85\% to 11.19\% but costs 18 clicks. Reinvesting 556~RUB at mobile-text CPC yields +14.8 expected clicks, netting 1{,}388 vs.\ 1{,}391 baseline ($-$3 clicks, $-$0.2\%) while lifting CTR to 11.31\%.

\begin{table}[ht]
\centering
\caption{Counterfactual format policy on historical May data.}
\label{tab:policy}
\begin{tabular}{lrrr}
\toprule
Scenario & Clicks & CTR (\%) & Spend (RUB) \\
\midrule
Observed (May) & 1{,}391 & 10.85 & 50{,}005 \\
Mobile graphical off & 1{,}373 & 11.19 & 49{,}450 \\
Off + budget to mobile text & 1{,}388 & 11.31 & $\sim$50{,}006 \\
\bottomrule
\end{tabular}
\end{table}

\subsection{Interpretation}

The format--device mismatch hypothesis is \emph{strongly supported}. Weakness is localized to graphical$\times$mobile, not graphical ads globally nor mobile traffic generally. Paired comparisons rule out the confound that graphical creatives were merely assigned to intrinsically weak keywords.

\section{Discussion and Limitations}

\paragraph{Comparing remedies.} Segment reallocation and format policy operate on different levers. Our simulations suggest format policy delivers a clearer CTR efficiency gain per ruble of identified waste (556~RUB at 3.32\% CTR), whereas group reallocation requires moving larger tranches for small click deltas. Practitioners should prioritize eliminating demonstrably weak creative--device pairs before fine-tuning keyword-group budgets.

\paragraph{Observational counterfactuals.} Simulations assume (i)~fixed CPC when scaling spend within a segment, (ii)~no auction feedback from policy changes, and (iii)~no seasonality beyond May 2024. Real A/B tests or geo experiments would be needed for causal confirmation.

\paragraph{Coverage gaps.} Targeting-category analysis (six query types; 83\% of budget in ``targeted'' queries at 12.5\% CTR) remains on the backlog. Tablet traffic is sparse and reported only diagnostically. Desktop text has only 14 impressions---unstable CTR.

\paragraph{External validity.} Results pertain to one regional bicycle retailer on Yandex Direct. Generalization to other verticals, platforms, or creative pipelines is unknown.

\section{Conclusion}

Structured hypothesis testing on a standard campaign export can separate \emph{where} budget is wasted from \emph{which mechanism} drives waste. For Centra Market in May 2024, both segment skew and mobile graphical underperformance are measurable; the latter is sharper, persistent across groups, and addressable with a simple format policy that nearly preserves click volume while improving CTR. Group-level reallocation is directionally correct but yields modest simulated gains, motivating larger or more targeted shifts. We release our analytical protocol as a template for small advertisers lacking custom ML infrastructure but seeking NeurIPS-style rigor in marketing post-mortems.

\begin{thebibliography}{9}
\bibitem{edelman2007internet}
B.~Edelman, M.~Ostrovsky, and M.~Schwarz.
\newblock Internet advertising and the generalized second-price auction: Selling billions of dollars worth of keywords.
\newblock \emph{American Economic Review}, 97(1):242--259, 2007.

\bibitem{ghosh2009ad}
A.~Ghosh, M.~Mahdian, S.~Estlahi, G.~Komijani, S.~Muthukrishnan, and B.~von Ahn.
\newblock To match or not to match: Economics of cookie matching in online advertising.
\newblock In \emph{Proceedings of EC}, 2014.

\bibitem{blake2015consumer}
T.~Blake, C.~M. Nosko, and S.~Tadelis.
\newblock Consumer heterogeneity and paid search effectiveness: A large-scale field experiment.
\newblock \emph{Econometrica}, 83(1):367--402, 2015.

\bibitem{ghose2013mobile}
A.~Ghose, S.~P.~Han, and J.~Park.
\newblock The economics of mobile applications.
\newblock \emph{Information Systems Research}, 24(4):613--631, 2013.

\bibitem{nuara2018combinatorial}
A.~Nuara, U.~Trov\`{o}, N.~Gatti, and M.~Restelli.
\newblock A combinatorial-bandit algorithm for the online joint bid/budget optimization of pay-per-click advertising campaigns.
\newblock In \emph{Proceedings of AAAI}, 2018.
\end{thebibliography}

\end{document}
